%%%%%%%%%%%%%%%%%%%%%%%%%%%%%%%%%%%%%%%%%%
\section{Proposed Approach}
\label{sec:proposed_approach}

\begin{figure}[htbp]
\centering
\includegraphics[width=\textwidth]{imagenes_paper/Diagrama_flujo.png}
\caption{Comprehensive workflow architecture of the proposed DTW/DDTW-driven adaptive rotational hybrid metaheuristic framework across heterogeneous problem domains, illustrating the closed-loop interaction between the multi-domain benchmark layer, the elastic stagnation monitoring engine, and the rotational solver alternation with elitist memory injection.}
\label{fig:framework_flowchart}
\end{figure}

The overall architecture and closed-loop operational workflow of the proposed collaborative optimization framework are presented in Figure~\ref{fig:framework_flowchart}. The framework integrates three interconnected functional layers: (i) a multi-domain problem layer covering discrete combinatorial optimization (MKP), continuous global benchmarks (IEEE CEC2022), and real-world engineering simulations (HRES2-H2 microgrid); (ii) an elastic DTW/DDTW stagnation monitoring engine that tracks sliding fitness trajectories ($\mathbf{X} \in \mathbb{R}^W$) against synthetic progress ($\mathbf{R}$) and stagnation ($\mathbf{C}$) baselines to compute distance metrics ($D_1, D_2, \Delta$) and confirm search stagnation via adaptive percentile thresholds and a patience streak filter ($S_t \ge 3$); and (iii) a collaborative metaheuristic execution layer that dynamically alternates active solvers between exploration ($\mathcal{P}_{pop}$) and exploitation ($\mathcal{P}_{traj}$) pools, seamlessly transferring search memory through an elitist solution injection protocol upon each validated solver switch.

\subsection{Overall Architecture of the DTW-Driven Adaptive Rotational Hybrid Framework}
\label{sec:overall_architecture}

To systematically resolve the critical challenges of premature convergence, progressive diversity loss, and limited cross-domain adaptability commonly encountered in standalone metaheuristic algorithms, this work proposes an Adaptive Rotational Hybrid Metaheuristic Framework governed by a Dynamic Time Warping (DTW) and Derivative Dynamic Time Warping (DDTW) Stagnation Monitor. Conventional hybrid schemes frequently switch between exploration and exploitation phases using rigid, heuristic criteria such as fixed iteration intervals or static stagnation counters which inevitably lead to either premature switching \cite{talbi2016combining}, (interrupting productive search trajectories) or delayed switching (wasting computational evaluations on stalled regions)\cite{akbulut2026artificial}. In contrast, the proposed framework employs temporal elastic alignment over the recent objective-function trajectory as a mathematically rigorous and dynamic indicator of search productivity  \cite{karimimamaghan2022machine}.

The theoretical foundation of the framework is rooted in the principle of algorithmic complementarity dictated by the No Free Lunch (NFL) theorem \cite{wolpert1997no}. Because no single search operator uniformly dominates across diverse landscape topologies, the proposed architecture organizes eight metaheuristic algorithms into two disjoint and specialized operational pools:
\begin{enumerate}
    \item \textbf{Population Pool ($\mathcal{P}_{pop}$):} Comprises population-based swarm and nature-inspired metaheuristics designed for global exploration and structural diversification across the search domain.
    \item \textbf{Trajectory Pool ($\mathcal{P}_{traj}$):} Comprises trajectory-based and local-search metaheuristics designed for fine grained exploitation and rapid intensification around promising regions.
\end{enumerate}

The execution workflow is structured into dynamic Epochs ($e = 1, 2, \dots, E$). In each epoch, an active solver $M^{(e)}$ randomly selected from the designated pool executes continuous optimization iterations on the objective function $f(\mathbf{x})$. Unlike static schedules, the duration of an epoch is intrinsically adaptive: the epoch terminates if and only if the DTW/DDTW Stagnation Monitor verifies that the recent convergence curve has entered an unpromising stagnation plateau. Upon triggering, the orchestrator freezes the active algorithm, extracts the global elite solution $\mathbf{x}_{best}^*$, alternates the pool type ($\mathcal{P}_{pop} \leftrightarrow \mathcal{P}_{traj}$), and executes the Seed Memory Injection Protocol to seamlessly initialize the successor solver without discarding accumulated evolutionary progress.



\subsection{Mathematical Formulation of the DTW/DDTW Stagnation Detection Engine}
\label{sec:dtw_formulation}

The stagnation detection engine operates on the sliding temporal sequence of best-so-far objective values recorded by the active solver across the most recent $W$ iterations:
\begin{equation}
\mathbf{X} = \left[ x_{t-W+1}, x_{t-W+2}, \dots, x_t \right] \in \mathbb{R}^W, \quad t \ge W
\label{eq:sliding_window}
\end{equation}

\subsubsection{Dynamic Construction of Baseline Reference Sequences}

To quantitatively evaluate whether $\mathbf{X}$ reflects productive convergence or an inactive plateau, the engine constructs two synthetic reference sequences of length $W$, both anchored at the initial value of the window $x_1 = x_{t-W+1}$:

\begin{enumerate}
    \item \textbf{Ideal Progress Ramp ($\mathbf{R} = [r_1, r_2, \dots, r_W]$):} Represents a steady, linear rate of improvement with minimum slope $s_{min}$:
    \begin{equation}
    r_k = x_1 + s_{min} \cdot (k - 1), \quad \forall k \in \{1, 2, \dots, W\}
    \label{eq:ramp_baseline}
    \end{equation}
    where $s_{min}$ is adaptively computed from the instantaneous dynamic range of the window:
    \begin{equation}
    s_{min} = 0.01 \cdot \frac{|x_W - x_1|}{W}
    \label{eq:adaptive_slope}
    \end{equation}

    \item \textbf{Constant Stagnation Plateau ($\mathbf{C} = [c_1, c_2, \dots, c_W]$):} Represents complete arrest of optimization progress:
    \begin{equation}
    c_k = x_1, \quad \forall k \in \{1, 2, \dots, W\}
    \label{eq:constant_baseline}
    \end{equation}
\end{enumerate}

\subsubsection{Dynamic Time Warping (DTW) with Sakoe--Chiba Band Constraint}

The elastic distance between the observed trajectory $\mathbf{X}$ and a target baseline sequence $\mathbf{Y} \in \{\mathbf{R}, \mathbf{C}\}$ is calculated by dynamic programming over the cumulative cost matrix $\mathbf{D} \in \mathbb{R}^{(W+1) \times (W+1)}$. The local cell distance is governed by the recurrence relation:
\begin{equation}
D(i, j) = |x_i - y_j| + \min \left\{ D(i-1, j),\, D(i, j-1),\, D(i-1, j-1) \right\}
\label{eq:dtw_recurrence}
\end{equation}
with boundary conditions $D(0, 0) = 0$ and $D(i, 0) = D(0, j) = +\infty$ for $i, j \ge 1$.

To prevent pathological warp alignments and enforce linear-time bounded computational complexity $\mathcal{O}(W \cdot w)$, the search space of the warping path is bounded by the Sakoe--Chiba band of width $w$:
\begin{equation}
|i - j| \le w, \quad w = \max(1, \lfloor \gamma \cdot W \rfloor), \quad \gamma = 0.10
\label{eq:sakoe_chiba}
\end{equation}
The final alignment metric corresponds to the terminal cell: $\text{DTW}(\mathbf{X}, \mathbf{Y}) = D(W, W)$.

\subsubsection{Derivative Dynamic Time Warping (DDTW) Formulation}

When optimization objective scales vary significantly or exhibit vertical shifts across distinct benchmark functions, standard DTW may remain sensitive to magnitude offsets. To ensure strict morphological shape-invariance, the framework implements Derivative Dynamic Time Warping (DDTW) \cite{keogh2001derivative}. The observed and reference series are mapped into their first-order finite differences:
\begin{equation}
\nabla x_k = x_k - x_{k-1}, \quad \text{with } \nabla x_1 = 0
\label{eq:derivative_transform}
\end{equation}
The DDTW distance is then computed by applying the DTW alignment procedure directly on the derivative sequences:
\begin{equation}
\text{DDTW}(\mathbf{X}, \mathbf{Y}) = \text{DTW}(\nabla \mathbf{X}, \nabla \mathbf{Y})
\label{eq:ddtw_metric}
\end{equation}

\subsubsection{Diagnostic Distance Metrics}

At each iteration $t \ge W$, the monitor extracts three primary diagnostic features:
\begin{align}
D_1 &= \text{DTW/DDTW}(\mathbf{X}, \mathbf{R}) \label{eq:d1_dist} \\
D_2 &= \text{DTW/DDTW}(\mathbf{X}, \mathbf{C}) \label{eq:d2_dist} \\
\Delta &= D_1 - D_2 \label{eq:delta_metric}
\end{align}
Under active convergence, $\mathbf{X}$ aligns closely with $\mathbf{R}$, yielding a low $D_1$, a large $D_2$, and a negative $\Delta < 0$. Conversely, stagnation causes $\mathbf{X}$ to collapse onto $\mathbf{C}$, resulting in $D_2 \to 0$, an elevated $D_1$, and $\Delta \gg 0$.

\subsection{Adaptive Thresholding via Moving Historical Percentiles}
\label{sec:adaptive_thresholding}

Static thresholding is inherently brittle when dealing with multi-domain benchmarks comprising disparate fitness landscapes (e.g., MKP with fitness in the order of $10^4$, CEC2022 in the order of $10^2$, and HRES2-H2 in sub-unit ranges). To provide scale-invariant portability, the framework maintains three cumulative memory buffers: $\mathcal{H}_{D_1}$, $\mathcal{H}_{D_2}$, and $\mathcal{H}_{\Delta}$, which store historical values of $D_1$, $D_2$, and $\Delta$ recorded throughout the current run.

Once a minimum buffer warm-up size is attained ($|\mathcal{H}| \ge 10$), the dynamic classification thresholds are updated at each iteration using empirical percentile functions:
\begin{align}
\theta_c &= \text{Percentile}\left(\mathcal{H}_{D_2}, P_{low}\right) \label{eq:theta_c} \\
\theta_r &= \text{Percentile}\left(\mathcal{H}_{D_1}, P_{high}\right) \label{eq:theta_r} \\
\theta_\Delta &= \text{Percentile}\left(\mathcal{H}_{\Delta}, P_{high}\right) \label{eq:theta_delta}
\end{align}
where $P_{low} = 30.0$ and $P_{high} = 70.0$.

The instantaneous triple stagnation condition $\text{Stagnant}(t) \in \{\text{True}, \text{False}\}$ is formally defined as:
\begin{equation}
\text{Stagnant}(t) \iff \left( N_{no\_improve} \ge K_{max} \right) \land \left( D_2 \le \theta_c \right) \land \left( (D_1 \ge \theta_r) \lor (\Delta \ge \theta_\Delta) \right)
\label{eq:stagnant_condition}
\end{equation}
where $N_{no\_improve}$ denotes the count of consecutive iterations without strict improvement in the local elite fitness.

To eliminate spurious triggers from short-term stochastic variance, a confirmation patience filter tracks the streak counter $S_t$:
\begin{equation}
S_t = \begin{cases}
S_{t-1} + 1, & \text{if } \text{Stagnant}(t) = \text{True} \\
0, & \text{if } \text{Stagnant}(t) = \text{False}
\end{cases}
\label{eq:streak_counter}
\end{equation}
The orchestrator triggers an algorithmic rotation event if and only if $S_t \ge P = 3$.




\subsection{Seed Memory Injection Protocol}

When the DTW monitor initiates solver rotation at epoch $e$, the orchestrator coordinates knowledge transfer via the following protocol:
\begin{enumerate}
    \item \textbf{Elite Extraction:} The global best solution $\mathbf{x}_{best}^* \in \mathbb{R}^D$ and its fitness $f(\mathbf{x}_{best}^*)$ are retrieved from global shared memory.
    \item \textbf{Injection into Population Solvers ($M \in \mathcal{P}_{pop}$):} The successor population of size $N_{pop}$ is initialized through a three-tier hybrid strategy:
    \begin{align}
    \mathbf{x}_1^{(0)} &= \mathbf{x}_{best}^* \label{eq:elite_inj} \\
    \mathbf{x}_i^{(0)} &= \mathbf{x}_{best}^* + \mathcal{N}\left(\mathbf{0}, \sigma^2 \mathbf{I}\right) \odot (\mathbf{U} - \mathbf{L}), \quad \forall i \in \left\{2, \dots, \left\lfloor \frac{N_{pop}}{2} \right\rfloor \right\} \label{eq:gauss_inj} \\
    \mathbf{x}_j^{(0)} &\sim \mathcal{U}(\mathbf{L}, \mathbf{U}), \quad \forall j \in \left\{\left\lfloor \frac{N_{pop}}{2} \right\rfloor + 1, \dots, N_{pop}\right\} \label{eq:unif_inj}
    \end{align}
    where $\sigma = 0.05$, and $[\mathbf{L}, \mathbf{U}]$ denotes the problem bounding box. This guarantees strong local exploitation around $\mathbf{x}_{best}^*$ while maintaining background exploratory dispersion.
    \item \textbf{Injection into Trajectory Solvers ($M \in \mathcal{P}_{traj}$):} The initial trajectory coordinate is assigned deterministically as $\mathbf{x}_0 = \mathbf{x}_{best}^*$, focusing immediate local intensification on the highest-quality basin.
\end{enumerate}
